(1) 未扰动本征能级为 $E_n^{(0)}=\hbar\omega(n+1/2)$，基态为零均值高斯，具有 $\langle x^2\rangle_0=\hbar/(2m\omega)$。记振子长度的方差参数 $\ell^2=\hbar/(2m\omega)$，位置算符可写为 $x=\sqrt{\hbar/(2m\omega)}(a+a^\dagger)$。
(2) 高斯基态的四阶矩由 Wick 配对或直接高斯积分给 $\langle x^4\rangle_0=3\langle x^2\rangle_0^2=3\hbar^2/(4m^2\omega^2)$。因此非简并微扰首阶能移 $\Delta E_0^{(1)}=\lambda\langle x^4\rangle_0=3\lambda\hbar^2/(4m^2\omega^2)>0$，正四次刚化势使基态能量增加。
(3) 用升降算符展开 $(a+a^\dagger)^4$ 并只保留在 $|n\rangle$ 的对角项，得到 $\langle n|x^4|n\rangle=3(2n^2+2n+1)\hbar^2/(4m^2\omega^2)$。故 $E_n\simeq\hbar\omega(n+1/2)+3\lambda(2n^2+2n+1)\hbar^2/(4m^2\omega^2)$，高能级间距因正四次势而增大。
(4) 取试探态 $\psi_s(x)=(\pi s^2)^{-1/4}e^{-x^2/(2s^2)}$，则 $\langle x^2\rangle=s^2/2$、$\langle p^2\rangle=\hbar^2/(2s^2)$、$\langle x^4\rangle=3s^4/4$。变分能量 $E(s)=\hbar^2/(4ms^2)+m\omega^2s^2/4+3\lambda s^4/4$。在 $\lambda=0$ 最小宽度 $s_0^2=\hbar/(m\omega)$，由于驻点性质，一阶能量修正只来自显式 $\lambda$ 项，等于 $3\lambda\hbar^2/(4m^2\omega^2)$。
